\documentclass[11pt]{article}
\usepackage[margin=1in]{geometry}
\usepackage{amsmath}
\usepackage{xcolor}
\usepackage{listings}
\usepackage{hyperref}
\hypersetup{colorlinks=true,linkcolor=blue,urlcolor=blue}

\definecolor{codebg}{rgb}{0.96,0.96,0.96}
\lstset{
  basicstyle=\ttfamily\footnotesize,
  backgroundcolor=\color{codebg},
  breaklines=true,
  frame=single,
  columns=fullflexible,
  keepspaces=true,
  showstringspaces=false,
  numbers=left, numberstyle=\tiny\color{gray},
}

\title{Proposal: reduce per-group allocation in \texttt{\_build\_protein\_rollup}}
\author{branch \texttt{perf/alloc-diagnostics}}
\date{2026-06-24}

\begin{document}
\maketitle

\section{Summary}
\texttt{\_build\_protein\_rollup} (in
\texttt{src/Routines/SearchDIA/SearchMethods/ProteinScoringSearch/protein\_inference\_pipeline.jl})
is called \textbf{once per protein group} inside two \texttt{groupby} loops. Each
call allocates \textbf{14 \texttt{Dict}s plus several row vectors}; over all groups
$\times$ files this is $\sim$0.6\,GB on the YeastStandard3M benchmark (the largest
single site in the Protein Scoring phase, which totals 1.98\,GB).

The fix is to (a) replace the 14 \texttt{Dict}s with \textbf{3 dedup \texttt{Dict}s +
reusable parallel arrays}, and (b) hoist that scratch out of the per-group function
so it is \textbf{allocated once and reused} across groups. Net per-group allocation
drops from ``14 Dicts + vectors'' to ``the returned output rows only''.

\textbf{Caveat (please note):} the current function emits its rows in
\texttt{Dict}-hash iteration order, and \texttt{pg\_score} is an
\emph{order-sensitive} \texttt{Float32} sum (\S\ref{sec:fragility}). The new design
iterates in deterministic first-seen order instead, so this is a \textbf{deliberate,
one-time protein-scoring rebaseline} (a few protein groups may shift, same class as
the approved iRT-refinement rebaseline). It also \emph{removes} the latent
order-fragility permanently. Verification plan in \S\ref{sec:verify}.

\section{Where it is called (per protein group)}
Loop 1 --- \texttt{protein\_inference\_pipeline.jl:664--666}:
\begin{lstlisting}[firstnumber=664]
for gdf in groupby(df, [:inferred_protein_group, :target, :entrap_id])
    quant_mask = _protein_rollup_quant_mask(gdf; q_value_threshold = q_value_threshold)
    rollup = _build_protein_rollup(gdf, quant_mask, prob_col)
\end{lstlisting}
Loop 2 --- \texttt{:1124--1126} (inside \texttt{combine(grouped) do gdf ... end}):
\begin{lstlisting}[firstnumber=1124]
protein_groups = combine(grouped) do gdf
    quant_mask = _protein_rollup_quant_mask(gdf; q_value_threshold = q_value_threshold)
    rollup = _build_protein_rollup(gdf, quant_mask, prob_col)
\end{lstlisting}
Both call once per \texttt{(inferred\_protein\_group, target, entrap\_id)} group, for
each run/file. Both loops are sequential (a plain \texttt{for}, and DataFrames
\texttt{combine}, which is single-threaded), so a single reusable scratch instance
per loop is safe.

\section{Current implementation (the 14 Dicts)}
\subsection*{Per-precursor level --- 7 Dicts (\texttt{:427--459})}
\begin{lstlisting}[firstnumber=427]
prob_by_precursor            = Dict{UInt32, Float32}()
best_peak_area_by_precursor  = Dict{UInt32, Float32}()
base_pep_id_by_precursor     = Dict{UInt32, UInt32}()
sequence_by_precursor        = Dict{UInt32, String}()
charge_by_precursor          = Dict{UInt32, UInt8}()
structural_mods_by_precursor = Dict{UInt32, String}()
isotopic_mods_by_precursor   = Dict{UInt32, String}()

@inbounds for i in eachindex(quant_mask)
    quant_mask[i] || continue
    precursor_idx = UInt32(gdf.precursor_idx[i])
    prob_val      = _sanitize_rollup_probability(gdf[!, prob_col][i])
    peak_area_val = Float32(gdf.peak_area[i])
    if haskey(prob_by_precursor, precursor_idx)
        prob_by_precursor[precursor_idx] = max(prob_by_precursor[precursor_idx], prob_val)
        if peak_area_val > best_peak_area_by_precursor[precursor_idx]
            best_peak_area_by_precursor[precursor_idx] = peak_area_val
        end
    else
        prob_by_precursor[precursor_idx]           = prob_val
        best_peak_area_by_precursor[precursor_idx] = peak_area_val
        base_pep_id_by_precursor[precursor_idx]    = UInt32(gdf.base_pep_id[i])
        sequence_by_precursor[precursor_idx]       = String(gdf.sequence[i])
        charge_by_precursor[precursor_idx]         = hasproperty(gdf, :charge) ? UInt8(gdf.charge[i]) : UInt8(0)
        structural_mods_by_precursor[precursor_idx] = ...   # String or ""
        isotopic_mods_by_precursor[precursor_idx]   = ...   # String or ""
    end
end
\end{lstlisting}
The output rows are then built by iterating \texttt{keys(prob\_by\_precursor)}
(\textbf{:463}) --- i.e.\ in \texttt{Dict}-hash order.

\subsection*{Per-modified-peptide level --- 4 Dicts (\texttt{:493--496})}
\begin{lstlisting}[firstnumber=493]
modified_log_none_sum   = Dict{Tuple{UInt32, String, String}, Float64}()
modified_best_peak_area = Dict{Tuple{UInt32, String, String}, Float32}()
modified_sequence       = Dict{Tuple{UInt32, String, String}, String}()
modified_precursor_count= Dict{Tuple{UInt32, String, String}, Int32}()
\end{lstlisting}
\subsection*{Per-peptide level --- 3 Dicts (\texttt{:529--531})}
\begin{lstlisting}[firstnumber=529]
peptide_log_none_sum   = Dict{String, Float64}()
peptide_best_peak_area = Dict{String, Float32}()
peptide_modified_count = Dict{String, Int32}()
\end{lstlisting}
Total: \textbf{7 + 4 + 3 = 14 Dicts allocated per call}. Measured contribution
(\texttt{Profile.Allocs}, sampled): $\sim$0.6\,GB across lines 438--540.

\section{The order fragility}\label{sec:fragility}
\texttt{pg\_score} is a \texttt{Float32} sum over \texttt{peptide\_rows}
(\textbf{:556}):
\begin{lstlisting}[firstnumber=556]
pg_score = isempty(peptide_rows) ? 0.0f0 : Float32(sum(row.score for row in peptide_rows))
\end{lstlisting}
\texttt{peptide\_rows} is built by iterating \texttt{keys(peptide\_log\_none\_sum)}
(\textbf{:545}), i.e.\ \texttt{Dict}-hash order. Floating-point summation is not
associative, so the iteration order changes \texttt{pg\_score} in the last bits.
\texttt{pg\_score} feeds protein-group FDR, so a last-bit change can flip a borderline
protein near the q-value threshold. \emph{This is why naive scratch reuse
(\texttt{empty!} changes a Dict's capacity, hence its iteration order) is not
byte-for-byte result-neutral} --- the same mechanism as the iRT \texttt{sizehint!}
issue. (\texttt{peptide\_list} at :558 is \texttt{sort!}'d so it is already
order-independent; \texttt{top\_pep\_peak\_area} is a \texttt{maximum}, also fine.)

\section{Why not merge the Dicts into one struct-valued Dict?}
A natural idea is \texttt{Dict\{UInt32, PrecAccum\}}. It does not help: the per-precursor
record holds \texttt{String} fields (\texttt{sequence}, \texttt{structural\_mods},
\texttt{isotopic\_mods}), so \texttt{PrecAccum} is not \texttt{isbits} and the \texttt{Dict}
would box one struct \emph{per precursor} --- worse than 7 scalar/String Dicts.

\section{Proposed change}\label{sec:proposal}
Use \textbf{one dedup \texttt{Dict} per level} (precursor\_idx / mod-key / sequence
$\rightarrow$ a row index) plus \textbf{parallel arrays} for the fields, pushed in
\textbf{first-seen order}. Iterate the arrays by position (deterministic,
capacity-independent). Hold all of this in a \texttt{ProteinRollupScratch} that the
caller allocates once and passes in; \texttt{\_reset!} it (\texttt{empty!} the 3
Dicts + the arrays, keeping their capacity) at the start of each call.

\subsection*{Scratch type (allocated once per loop, reused)}
\begin{lstlisting}
mutable struct ProteinRollupScratch
    # precursor level (dedup by precursor_idx; parallel arrays, first-seen order)
    prec_pos    :: Dict{UInt32, Int}
    prec_id     :: Vector{UInt32}
    prec_prob   :: Vector{Float32}
    prec_peak   :: Vector{Float32}
    prec_basepep:: Vector{UInt32}
    prec_seq    :: Vector{String}
    prec_charge :: Vector{UInt8}
    prec_smods  :: Vector{String}
    prec_imods  :: Vector{String}
    # modified-peptide level (dedup by (base_pep_id, structural_mods, isotopic_mods))
    mod_pos     :: Dict{Tuple{UInt32,String,String}, Int}
    mod_key     :: Vector{Tuple{UInt32,String,String}}
    mod_logsum  :: Vector{Float64}
    mod_peak    :: Vector{Float32}
    mod_seq     :: Vector{String}
    mod_count   :: Vector{Int32}
    # peptide level (dedup by sequence)
    pep_pos     :: Dict{String, Int}
    pep_seq     :: Vector{String}
    pep_logsum  :: Vector{Float64}
    pep_peak    :: Vector{Float32}
    pep_count   :: Vector{Int32}
end
\end{lstlisting}

\subsection*{Per-precursor accumulation (replaces :427--459)}
\begin{lstlisting}
_reset!(s)   # empty! the 3 Dicts and all arrays (capacity retained for reuse)
@inbounds for i in eachindex(quant_mask)
    quant_mask[i] || continue
    pid       = UInt32(gdf.precursor_idx[i])
    prob_val  = _sanitize_rollup_probability(gdf[!, prob_col][i])
    peak_val  = Float32(gdf.peak_area[i])
    pos = get(s.prec_pos, pid, 0)
    if pos == 0                              # first time we see this precursor
        push!(s.prec_id, pid)
        push!(s.prec_prob, prob_val)
        push!(s.prec_peak, peak_val)
        push!(s.prec_basepep, UInt32(gdf.base_pep_id[i]))
        push!(s.prec_seq, String(gdf.sequence[i]))
        push!(s.prec_charge, hasproperty(gdf, :charge) ? UInt8(gdf.charge[i]) : UInt8(0))
        push!(s.prec_smods, _mods_or_empty(gdf, :structural_mods, i))
        push!(s.prec_imods, _mods_or_empty(gdf, :isotopic_mods, i))
        s.prec_pos[pid] = length(s.prec_id) # row index in the parallel arrays
    else                                     # seen before: update aggregates only
        s.prec_prob[pos] = max(s.prec_prob[pos], prob_val)
        if peak_val > s.prec_peak[pos]; s.prec_peak[pos] = peak_val; end
    end
end
# build precursor_rows by iterating 1:length(s.prec_id)  (first-seen order)
\end{lstlisting}
The modified-peptide and peptide levels follow the identical pattern (one dedup
\texttt{Dict} + parallel arrays each). The returned
\texttt{precursor\_rows / modified\_peptide\_rows / peptide\_rows} are still built
(they are the output the callers consume) --- only the \emph{intermediate} 14 Dicts
become 3 reused dedup Dicts + reused arrays.

\subsection*{Caller change (allocate scratch once)}
\begin{lstlisting}
scratch = ProteinRollupScratch()                     # once, before the loop
for gdf in groupby(df, [...])
    quant_mask = _protein_rollup_quant_mask(gdf; ...)
    rollup = _build_protein_rollup(gdf, quant_mask, prob_col, scratch)   # 4-arg
end
\end{lstlisting}
A 3-arg method \texttt{\_build\_protein\_rollup(gdf, quant\_mask, prob\_col)} is kept
(allocates a fresh scratch then calls the 4-arg form) so the existing signature and
any tests keep working.

\section{Allocation impact (expected)}
Per group: \textbf{14 Dict allocations $\rightarrow$ 0} (3 dedup Dicts + 14 arrays
all reused from the scratch). The only remaining per-group allocations are the
returned output row vectors (necessary). Expected saving on YeastStandard3M:
$\sim$0.3--0.5\,GB of the 0.6\,GB (to be confirmed by the
\texttt{@alloc\_bucket}/\texttt{Profile.Allocs} re-run).

\section{Correctness, rebaseline, and verification}\label{sec:verify}
\begin{itemize}
  \item \textbf{Aggregates unchanged:} \texttt{max}/\texttt{sum}/\texttt{count} per
        precursor/peptide are identical regardless of iteration order; only the
        \emph{order} in which rows are emitted changes (Dict-hash $\rightarrow$
        first-seen).
  \item \textbf{Deliberate rebaseline:} because \texttt{pg\_score} sums in row order,
        the new first-seen order changes \texttt{pg\_score} in the last bits, so a
        few borderline protein groups may cross the FDR threshold. This is a one-time
        rebaseline (and removes the latent order-fragility for good --- the result
        becomes a deterministic function of the input, not of Dict capacity).
  \item \textbf{Verification:} fresh-process determinism (the new protein-group /
        precursor counts must be \emph{reproducible} across runs); confirm the change
        vs the current baseline is small ($\lesssim$1\,\%, like the iRT rebaseline)
        --- a large change would signal a bug. Re-run \texttt{@alloc\_bucket} /
        \texttt{Profile.Allocs} to confirm the allocation drop and no Protein Scoring
        time regression. Existing ProteinScoringSearch unit tests must pass.
\end{itemize}

\section{Optional follow-on (not in this change)}
The returned \texttt{*\_rows} \texttt{Vector\{NamedTuple\}} are still allocated per
group. If the callers consumed the scratch parallel arrays directly, those could be
avoided too --- but that touches the caller contract and is left out to keep this
change focused.

\end{document}
